% GENERATED FILE. Edit canonical lessons, then run npm run generate:books.
% canonical-source-hash: fnv1a64:ebb6a9ef76bced1c
% canonical-lessons: FR-C53-douleur, FR-C53-fievre, FR-C53-toux, FR-C53-medicament, FR-C53-medecin

\chapter{Pain, Fever, Cough, Medicine, Doctor}
\label{ch:fr-pain-fever-cough-medicine-doctor}
\hlchaptermodality{53}

\begin{chapteropening}
\textbf{By the end of this chapter:} \emph{I can name pain, a fever, a cough, a medicine, a doctor.}

Complete the last lesson of chapter 53: I can name pain, a fever, a cough, a medicine, a doctor.
\end{chapteropening}

\section[la douleur]{la douleur — pain}

\label{lesson:FR-C53-douleur}



\begin{quote}
Before the new one: say the French for an elbow, then the French for a wrist.
\end{quote}

\subsection*{You'll want to know: la douleur}
\textbf{la douleur} — \textquotedblleft{}pain\textquotedblright{}.

\begin{cousinweb}
\textbf{douleur} $\leftarrow$ Old French \emph{dolor} $\leftarrow$ Latin \emph{dolor}, \textquotedblleft{}pain, grief\textquotedblright{} — inherited. (EtymDB 2.1)
\end{cousinweb}

\begin{grammarlens}[title={What you've built}]
One more word for this chapter.
\end{grammarlens}

\subsection*{Your turn}
\begin{itemize}
\raggedright
  \item \emph{Say it:} \emph{la douleur}
  \item \emph{Say it:} \emph{la douleur}, once more
  \item \emph{Say it:} say \emph{la douleur}
  \item \emph{Recall:} say the French for an elbow, then the French for a wrist, then say \emph{la douleur} again
\end{itemize}

\begin{tcolorbox}[breakable,colback=teal!4,colframe=teal!35!black,title={Before you move on}]
What does la douleur mean? (\textquotedblleft{}Pain\textquotedblright{}.) Say it once more.
\end{tcolorbox}
\section[la fièvre]{la fièvre — a fever}

\label{lesson:FR-C53-fievre}



\begin{quote}
Before the new one: say the French for a wrist, then the French for pain.
\end{quote}

\subsection*{You'll want to know: la fièvre}
\textbf{la fièvre} — \textquotedblleft{}a fever\textquotedblright{}.

\begin{cousinweb}
\textbf{fièvre} $\leftarrow$ Latin \emph{febris}, \textquotedblleft{}fever\textquotedblright{} (EtymDB 2.1) — inherited: the stressed \emph{e} broke into \emph{ie}, and \emph{b} became \emph{v}.
\end{cousinweb}

\begin{grammarlens}[title={What you've built}]
One more word for this chapter.
\end{grammarlens}

\subsection*{Your turn}
\begin{itemize}
\raggedright
  \item \emph{Say it:} \emph{la fièvre}
  \item \emph{Say it:} \emph{la fièvre}, once more
  \item \emph{Say it:} say \emph{la fièvre}
  \item \emph{Recall:} say the French for a wrist, then the French for pain, then say \emph{la fièvre} again
\end{itemize}

\begin{tcolorbox}[breakable,colback=teal!4,colframe=teal!35!black,title={Before you move on}]
What does la fièvre mean? (\textquotedblleft{}A fever\textquotedblright{}.) Say it once more.
\end{tcolorbox}
\section[la toux]{la toux — a cough}

\label{lesson:FR-C53-toux}



\begin{quote}
Before the new one: say the French for pain, then the French for a fever.
\end{quote}

\subsection*{You'll want to know: la toux}
\textbf{la toux} — \textquotedblleft{}a cough\textquotedblright{}.

\begin{cousinweb}
\textbf{toux} $\leftarrow$ Old French \emph{tous} (EtymDB 2.1) $\leftarrow$ Latin \emph{tussis}, \textquotedblleft{}a cough\textquotedblright{} (Lewis \& Short) — inherited.
\end{cousinweb}

\begin{grammarlens}[title={What you've built}]
One more word for this chapter.
\end{grammarlens}

\subsection*{Your turn}
\begin{itemize}
\raggedright
  \item \emph{Say it:} \emph{la toux}
  \item \emph{Say it:} \emph{la toux}, once more
  \item \emph{Say it:} say \emph{la toux}
  \item \emph{Recall:} say the French for pain, then the French for a fever, then say \emph{la toux} again
\end{itemize}

\begin{tcolorbox}[breakable,colback=teal!4,colframe=teal!35!black,title={Before you move on}]
What does la toux mean? (\textquotedblleft{}A cough\textquotedblright{}.) Say it once more.
\end{tcolorbox}
\section[le médicament]{le médicament — a medicine}

\label{lesson:FR-C53-medicament}



\begin{quote}
Before the new one: say the French for a fever, then the French for a cough.
\end{quote}

\subsection*{You'll want to know: le médicament}
\textbf{le médicament} — \textquotedblleft{}a medicine\textquotedblright{}.

\begin{cousinweb}
\textbf{médicament} $\leftarrow$ Latin \emph{medicamentum}, \textquotedblleft{}a remedy\textquotedblright{}, from \emph{medicari}, \textquotedblleft{}to heal\textquotedblright{} — a learned borrowing. (EtymDB 2.1)
\end{cousinweb}

\begin{grammarlens}[title={What you've built}]
One more word for this chapter.
\end{grammarlens}

\subsection*{Your turn}
\begin{itemize}
\raggedright
  \item \emph{Say it:} \emph{le médicament}
  \item \emph{Say it:} \emph{le médicament}, once more
  \item \emph{Say it:} say \emph{le médicament}
  \item \emph{Recall:} say the French for a fever, then the French for a cough, then say \emph{le médicament} again
\end{itemize}

\begin{tcolorbox}[breakable,colback=teal!4,colframe=teal!35!black,title={Before you move on}]
What does le médicament mean? (\textquotedblleft{}A medicine\textquotedblright{}.) Say it once more.
\end{tcolorbox}
\section[le médecin]{le médecin — a doctor}

\label{lesson:FR-C53-medecin}



\begin{quote}
Before the new one: say the French for a cough, then the French for a medicine.
\end{quote}

\subsection*{You'll want to know: le médecin}
\textbf{le médecin} — \textquotedblleft{}a doctor\textquotedblright{}.

\begin{cousinweb}
\textbf{médecin} — formed in French from the word for \textquotedblleft{}medicine\textquotedblright{}, which $\leftarrow$ Latin \emph{medicina}, \textquotedblleft{}the art of healing\textquotedblright{}, from \emph{medicus}, \textquotedblleft{}doctor\textquotedblright{} (Lewis \& Short).
\end{cousinweb}

\begin{grammarlens}[title={What you've built}]
One more word for this chapter.
\end{grammarlens}

\subsection*{Your turn}
\begin{itemize}
\raggedright
  \item \emph{Say it:} \emph{le médecin}
  \item \emph{Say it:} \emph{le médecin}, once more
  \item \emph{Say it:} say \emph{le médecin}
  \item \emph{Recall:} say the French for a cough, then the French for a medicine, then say \emph{le médecin} again
\end{itemize}

\begin{tcolorbox}[breakable,colback=teal!4,colframe=teal!35!black,title={Before you move on}]
What does le médecin mean? (\textquotedblleft{}A doctor\textquotedblright{}.) Say it once more.
\end{tcolorbox}
